% 三角不等式与反向三角不等式：(a) x=(0,0), y=(2.2,1.5), z=(3.6,0)，直达与绕行；(b) 同一三角形，|d(x,z)-d(y,z)|<=d(x,y)。
\begin{tikzpicture}[
  pt/.style={circle,fill=FigureStroke,inner sep=0pt,minimum size=4pt},
  direct/.style={draw=FigureBlue,line width=1.1pt},
  detour/.style={draw=FigureMuted,line width=.8pt,dashed},
  note/.style={font=\FigureNoteFont\fontsize{7.5}{9}\selectfont,text=FigureMuted,align=center}]
\begin{scope}
  \coordinate (x) at (0,0); \coordinate (y) at (2.2,1.5); \coordinate (z) at (3.6,0);
  \draw[direct] (x)--(z) node[midway,below=1pt,text=FigureInk]{$d(x,z)$};
  \draw[detour] (x)--(y) node[midway,above left=-1pt,text=FigureInk]{$d(x,y)$};
  \draw[detour] (y)--(z) node[midway,above right=-1pt,text=FigureInk]{$d(y,z)$};
  \node[pt,label=left:$x$] at (x){}; \node[pt,label=above:$y$] at (y){}; \node[pt,label=right:$z$] at (z){};
  \node[note] at (1.8,-.9) {直达不长于绕行};
  \node at (1.8,-1.5) {(a) $d(x,z)\leq d(x,y)+d(y,z)$};
\end{scope}
\begin{scope}[xshift=5.2cm]
  \coordinate (x) at (0,0); \coordinate (y) at (2.2,1.5); \coordinate (z) at (3.6,0);
  \draw[direct] (x)--(z);
  \draw[draw=FigureStroke,line width=.8pt] (y)--(z);
  \draw[detour] (x)--(y) node[midway,above left=-1pt,text=FigureInk]{$d(x,y)$};
  % 以 z 为圆心、半径 d(y,z) 的圆弧与 xz 交于 w，xw 的长度即两距离之差
  \pgfmathsetmacro{\ryz}{sqrt(1.4*1.4+1.5*1.5)}
  \draw[FigureGrid,line width=.6pt] (z) ++(180:\ryz) arc (180:133:\ryz);
  \coordinate (w) at ($(z)+(-\ryz,0)$);
  \draw[draw=FigureGreen,line width=2pt] (x)--(w);
  \node[pt,label=left:$x$] at (x){}; \node[pt,label=above:$y$] at (y){}; \node[pt,label=right:$z$] at (z){};
  \node[note,text=FigureGreen] at ($(x)!.5!(w)+(0,-.35)$) {差};
  \node[note] at (1.8,-.9) {到$z$的距离之差不超过$d(x,y)$};
  \node at (1.8,-1.5) {(b) $|d(x,z)-d(y,z)|\leq d(x,y)$};
\end{scope}
\end{tikzpicture}
